\documentclass[10pt]{article}
\usepackage[T1]{fontenc}
\usepackage[utf8]{inputenc}
\usepackage[french]{babel}
\usepackage{amsmath,amssymb,mathrsfs,mathpazo}
\usepackage{geometry}
\geometry{a4paper,margin=2cm}
\usepackage{graphicx,array,tabularx,tabularray,colortbl,multirow,multicol,booktabs,enumitem,pifont,ulem,fancybox,diagbox,cancel}
\usepackage{pstricks,pstricks-add,pst-plot,pst-tree,pst-node,pst-text,pst-func,pst-math,pst-bspline,pst-3dplot,pst-eucl,pst-solides3d}
\let\pscalculate\undefined
\usepackage{pst-bezier}
\usepackage{tikz,pgfplots,tkz-tab}
\usetikzlibrary{arrows,arrows.meta,patterns,calc,positioning,trees,decorations.markings,shapes.geometric,shapes.misc,shapes.arrows,intersections,decorations.pathmorphing,decorations.pathreplacing,backgrounds,shadows,babel}
\pgfplotsset{compat=1.18}
\usepgfplotslibrary{fillbetween}
\usepackage[np]{numprint}
\usepackage{listings,makecell,pas-tableur,twoopt}
\let\barre\undefined
\usepackage[pstricks]{tablvar}
\definecolor{darkgreen}{RGB}{0,100,0}
\providecommand{\red}{\color{red}}
\providecommand{\blue}{\color{blue}}
\providecommand{\green}{\color{green}}
\providecommand{\black}{\color{black}}
\providecommand{\white}{\color{white}}
\providecommand{\gray}{\color{gray}}
\providecommand{\cyan}{\color{cyan}}
\providecommand{\magenta}{\color{magenta}}
\providecommand{\yellow}{\color{yellow}}
\providecommand{\violet}{\color{violet}}
\providecommand{\orange}{\color{orange}}
\providecommand{\euro}{\texteuro}
\providecommand{\EUR}{\texteuro}
\providecommand{\ieme}{\textsuperscript{e}}
\providecommand{\degres}{\ensuremath{^\circ}}
\providecommand{\degree}{\ensuremath{^\circ}}
\providecommand{\vv}[1]{\overrightarrow{#1}}
\providecommand{\Coord}[1]{#1}
\providecommand{\decosix}{}
\providecommand{\starredbullet}{\ensuremath{\bullet}}
\pagestyle{empty}
\setlength{\parindent}{0pt}
\usepackage{adjustbox}\DeclareUnicodeCharacter{2236}{:}\DeclareUnicodeCharacter{2190}{\ensuremath{\leftarrow}}
\providecommand{\R}{\mathbb{R}}
\providecommand{\N}{\mathbb{N}}
\providecommand{\Z}{\mathbb{Z}}
\providecommand{\Q}{\mathbb{Q}}
\providecommand{\D}{\mathbb{D}}
\providecommand{\C}{\mathbb{C}}
\providecommand{\e}{{\mathrm{e}}}
\providecommand{\vect}[1]{\overrightarrow{#1}}
\providecommand{\vectt}[1]{\overrightarrow{\mathrm{#1}}}
\providecommand{\vvt}[1]{\overrightarrow{\mathrm{#1}}}
\providecommand{\np}[2][]{{#2~#1}}
\providecommand{\vv}[1]{{\overrightarrow{#1}}}
\renewcommand{\d}{{\mathrm{d}}}
\providecommand{\degres}{{^\circ}}
\providecommand{\degre}{{^\circ}}
\providecommand{\blue}{\color{blue}}
\providecommand{\red}{\color{red}}
\providecommand{\enskip}{\quad}
\providecommand{\llbracket}{\lbrack\!\lbrack}
\providecommand{\rrbracket}{\rbrack\!\rbrack}
\providecommand{\euro}{{\text{€}}}
\providecommand{\no}{{\text{n°}}}
\providecommand{\interval}[2]{{[#1;#2]}}
\providecommand{\intervFO}[2]{{[#1;#2[}}
\providecommand{\dotfill}{\ldots}

\providecommand{\ds}{\displaystyle}
\providecommand{\textding}[1]{#1}
\providecommand{\Oij}{\(\mathrm{O};\vec{\imath},\vec{\jmath}\)}
\providecommand{\Oijk}{\(\mathrm{O};\vec{\imath},\vec{\jmath},\vec{k}\)}
\providecommand{\Ouv}{\(\mathrm{O};\vec{u},\vec{v}\)}
\providecommand{\intervFF}[2]{\ensuremath{[#1;#2]}}
\renewcommand{\euro}{\ifmmode\text{\texteuro}\else\texteuro\fi}
\DeclareUnicodeCharacter{00B0}{\ifmmode^{\circ}\else\textdegree\fi}
\DeclareUnicodeCharacter{20AC}{\ifmmode\text{\texteuro}\else\texteuro\fi}
% Énoncés : Éducation nationale. Transcription et corrigés : APMEP, Denis Vergès et contributeurs. Découpage et classement : CoopMaths, Sébastien Lozano, Éric Elter et contributeurs. Adaptation de la présentation : Maths Explorer.
% CC BY-SA : https://coopmaths.fr/www/about/
% https://coopmaths.fr/alea/static/bac/2024/tex/bac_2024_05_sujet1_ameriquenord_4.tex
\begin{document}
\noindent\textit{APMEP / CoopMaths. CC BY-SA. Présentation adaptée par Maths Explorer.}\par\medskip
\medskip

Pour tout entier naturel $n$, on considère les intégrales suivantes :

\[I_n = \displaystyle\int_0^{\pi} \text{e}^{- nx}\sin (x)\:\text{d}x, \quad  J_n = \displaystyle\int_0^{\pi} \text{e}^{-nx}\cos (x)\:\text{d}x.\]

\medskip

\begin{enumerate}
\item Calculer $I_0$.
\item 
	\begin{enumerate}
		\item Justifier que, pour tout entier naturel $n$ , on a $I_n \geqslant 0$.
		\item Montrer que, pour tout entier naturel $n$ , on a $I_{n+1} - I_n \leqslant 0$.
		\item Déduire des deux questions précédentes que la suite $\left(I_n\right)$ converge.
	\end{enumerate}
\item
	\begin{enumerate}
		\item Montrer que, pour tout entier naturel $n$ , on a : 
		\[I_n \leqslant \displaystyle\int_0^{\pi} \text{e}^{- nx}\:\text{d}x.\]
		

		\item Montrer que, pour tout entier naturel $n \geqslant 1$, on a : 
		\[\displaystyle\int_0^{\pi} \text{e}^{- nx}\:\text{d}x = \dfrac{1 - \e^{-n \pi}}{n}.\]
		
		\item Déduire des deux questions précédentes la limite de la suite $\left(I_n\right)$.
	\end{enumerate}
\item
	\begin{enumerate}
		\item En intégrant par parties l'intégrale $I_n$ de deux façons différentes, établir les deux relations suivantes, pour tout entier naturel $n \geqslant 1$ :

\begin{center}$I_n = 1 + \e^{-n\pi} - nJ_n$ \quad et\quad  $I_n = \dfrac 1n J_n$\end{center}

		\item En déduire que, pour tout entier naturel $n \geqslant 1$, on a 
		\[I_n = \dfrac{1 + \e^{-n \pi}}{n^2 + 1}\]
	\end{enumerate}
\item On souhaite obtenir le rang $n$ à partir duquel la suite $\left(I_n\right)$ devient inférieure à 0,1.

Recopier et compléter la cinquième ligne du script Python ci-dessous avec la
commande appropriée.
\begin{center}
\begin{tabular}{|l l|}\hline
1 &\textbf{from} math \textbf{import *}\\
2 &def seuil() :\\
3 &\quad n = 0\\
4 &\quad I = 2\\
5 &\quad \ldots\\
6 &\quad \qquad n=n+1\\
7 &\quad \qquad I=(1+exp(-n*pi))/(n*n+1)\\
8 & \textbf{return} n\\ \hline
\end{tabular}
\end{center}

\end{enumerate}

\end{document}
